%
% PhD thesis style
% 
% improved by Vincent X. Gong
% with love
%
% October, 2019
% http://www.gong.im
% http://https://www.gong.im/phd_thesis_latex_trail_template/
%
% originally created by Rudy Negenborn
%
% based on trail thesis style for LaTeX
%
% February 2008
%
% further info: rudy@negenborn.net
% and: http://www.negenborn.net/mampc/
%

% How to get started? For a start:
% * read through this file
% * edit 0-misc/defs.tex
% * edit 0-misc/title.tex
% * latex thesis.tex
% and there you have your first version

\documentclass[12pt,hidelinks]{trailthesis}
%\documentclass[a4paper,12pt]{book}


% include the draft option to include tags in the .dvi showing were
% overfull hboxes are located, e.g.,
% \documentclass[10pt,draft]{trailthesis}

% include the hidelinks option to remove the red box in the table of content

%%%%%%%%%%%%%%%%%%%%%%%%%%%%%%%%%%%%%%%%%%%%%%%%%%%%%%%%%%%%%%%%

% define the base directory in which your thesis files are located
%
%
% NOTE WINDOWS USERS: the '\' in the Windows directory name should
%                     be replaced by '/', so:
%                     C:/thesis/, and not C:\thesis\
%
% NOTE WINDOWS USERS: make sure that you thesis files are in
%                     a directory *without* spaces in the name,
%                     so: C:/thesis, and not C:/My Document/My Thesis
%
% NOTE: make sure that the directory name ends with a '/'
%
%\newcommand{\basedir}{/Users/vgong/Documents/workspaces/phd_thesis/thesis_trail_vg/}

%\usepackage[xetex]{hyperref}
%\usepackage{hyperref}
%%%%%%%%%%%%%%%%%%%%%%%%%%%%%%%%%%%%%%%%%%%%%%%%%%%%%%%%%%%%%%%

% include following package to be able to have line numbers in from of
% all text of thesis. this can be useful to include in the draft
% version of your thesis, such that the committee members can easily
% point out where they have remarks/suggestions/corrections.

\usepackage[left,modulo]{lineno}

%%%%%%%%%%%%%%%%%%%%%%%%%%%%%%%%%%%%%%%%%%%%%%%%%%%%%%%%%%%%%%%%

\newcommand\dedication[1]{%
	\thispagestyle{empty}%
	\vspace*{\fill}%
	\begin{flushright}%
		#1%
	\end{flushright}%
%	\vspace*{\fill}%
%	\cleardoublepage%
}
\newcommand\epigraph[3][2\baselineskip]{%
	\begin{flushright}%
		{\rmfamily\itshape #2}%
		\vskip 0.5\baselineskip%
		#3%
	\end{flushright}%
	\vspace{#1}%
}



% include base definitions used in this thesis
%%%%%%%%%%%%%%%%%%%%%%%%%%%%%%%%%%%%%%%%%%%%%%%%%%%%%%%%%%%%%

% define some commands for handling images in subdirectories

\newcommand{\fname}[1]{\basedir#1}

% command for dealing with images
\newcommand{\imagepath}[1]{\fname{1-intro/images/#1}}


% define note command

\setlength{\marginparwidth}{30mm}

\newcommand{\mynote}[1]{%
\marginpar{\footnotesize{#1}}%
}%

% uncomment to remove remarks.
%\renewcommand{\mynote}[1]{} 

%%%%%%%%%%%%%%%%%%%%%%%%%%%%%%%%%%%%%%%%%%%%%%%%%%%%%%%%%%%%%

% define psfrag centering command
\newcommand{\psfragc}[2]{%
\psfrag{#1}[t][t]{\color[rgb]{0,0,0}%
\setlength{\tabcolsep}{0pt}%
\begin{tabular}{c}#2\end{tabular}}%
}


%%%%%%%%%%%%%%%%%%%%%%%%%%%%%%%%%%%%%%%%%%%%%%%%%%%%%%%%%%%%%
% define the bib style

\usepackage[authoryear]{natbib}
%\usepackage[]{natbib}

%%%%%%%%%%%%%%%%%%%%%%%%%%%%%%%%%%%%%%%%%%%%%%%%%%%%%%%%%%%%%

% adjust maximum number of columns that can be used in a matrix environment

% add matrix cols
%\setcounter{MaxMatrixCols}{30}



%%% Local Variables: 
%%% mode: latex
%%% TeX-master: "../2007-thesis"
%%% End: 

%%%%%%%%%%%%%%%%%%%%%%%%%%%%%%%%%%%%%%%%%%%%%%%%%%%%%%%%%%%%%%

% define some commands for handling images in subdirectories

\newcommand{\fname}[1]{\basedir#1}

% command for dealing with images
\newcommand{\imagepath}[1]{\fname{1-intro/images/#1}}


% define note command

\setlength{\marginparwidth}{30mm}

\newcommand{\mynote}[1]{%
\marginpar{\footnotesize{#1}}%
}%

% uncomment to remove remarks.
%\renewcommand{\mynote}[1]{} 

%%%%%%%%%%%%%%%%%%%%%%%%%%%%%%%%%%%%%%%%%%%%%%%%%%%%%%%%%%%%%

% define psfrag centering command
\newcommand{\psfragc}[2]{%
\psfrag{#1}[t][t]{\color[rgb]{0,0,0}%
\setlength{\tabcolsep}{0pt}%
\begin{tabular}{c}#2\end{tabular}}%
}


%%%%%%%%%%%%%%%%%%%%%%%%%%%%%%%%%%%%%%%%%%%%%%%%%%%%%%%%%%%%%
% define the bib style

\usepackage[authoryear]{natbib}
%\usepackage[]{natbib}

%%%%%%%%%%%%%%%%%%%%%%%%%%%%%%%%%%%%%%%%%%%%%%%%%%%%%%%%%%%%%

% adjust maximum number of columns that can be used in a matrix environment

% add matrix cols
%\setcounter{MaxMatrixCols}{30}



%%% Local Variables: 
%%% mode: latex
%%% TeX-master: "../2007-thesis"
%%% End: 


%%%%%%%%%%%%%%%%%%%%%%%%%%%%%%%%%%%%%%%%%%%%%%%%%%%%%%%%%%%%%%%%

% include other definitions that you have defined yourself
%%%%%%%%%%%%%%%%%%%%%%%%%%%%%%%%%%%%%%%%%%%%%%%%%%%%%%%%%%%%%

% include here your own macros


%%% Local Variables: 
%%% mode: latex
%%% TeX-master: "../thesis"
%%% End: 

\usepackage{multirow}
\usepackage{booktabs}

% Pages with this attribute will be displayed in landscape orientation by conforming PDF viewers
\usepackage{pdflscape}
\usepackage{hhline}
\usepackage{enumitem}
%\usepackage{subcaption}
\usepackage{float}
\usepackage{threeparttable}
\usepackage{tikz} % FvW: to make very nice looking graphs
\usepackage{CJKutf8} % for insert Chinese
\usepackage[utf8]{inputenc}
\usepackage{wrapfig}
\usepackage{etex}



\begin{document}

%%%%%%%%%%%%%%%%%%%%%%%%%%%%%%%%%%%%%%%%%%%%%%%%%%%%%%%%%%%%%%%%

\frontmatter

%% define relative path to images for this .tex file
\graphicspath{{./0-misc/images/}}


%%%%%%%%%%%%%%%%%%%%%%%%%%%%%%%%%%%%%%%%%%%%%%%%%%%%%%%%%%%%%%%%%%%%
% define below the title of your thesis, the subtitle, and your own name

\newcommand{\thetitle}{Using Social Media to Characterise Crowds in City Events for Crowd Management} 
\newcommand{\thesubtitle}{} 
\newcommand{\theauthor}{Vincent\ X.\ GONG}


%%%%%%%%%%%%%%%%%%%%%%%%%%%%%%%%%%%%%%%%%%%%%%%%%%%%%%%%%%%%%%%%%%%%
%% First the initial page with only the title and the author name
%%%%%%%%%%%%%%%%%%%%%%%%%%%%%%%%%%%%%%%%%%%%%%%%%%%%%%%%%%%%%%%%%%%%
\pagestyle{fancy} %% start page headers and footers
\thispagestyle{empty}
\vspace*{30mm}

\begin{center}
    \textbf{\huge \thetitle}\\[\baselineskip]\textbf{\LARGE \thesubtitle}
\end{center}


\vspace{40mm}

\begin{center}
    \Large\theauthor
\end{center}

%%%%%%%%%%%%%%%%%%%%%%%%%%%%%%%%%%%%%%%%%%%%%%%%%%%%%%%%%%%%%%%%%%%%
%% the back (page 2)
%%%%%%%%%%%%%%%%%%%%%%%%%%%%%%%%%%%%%%%%%%%%%%%%%%%%%%%%%%%%%%%%%%%%
\newpage

\thispagestyle{empty}

\vspace*{\fill}
% possibly include your logos and make acknowledgements to funding

\noindent \includegraphics[width=7em]{erc.png}\hspace*{0.75cm}
\includegraphics[width=7em]{allegro.png}\hspace*{0.75cm}
\includegraphics[width=12.5em]{ams.png}\hspace*{0.75cm}\\
%\noindent \includegraphics[width=21em]{logo_allegro}\\

\noindent The research leading to this dissertation has received funding from the European Research Council under the European Union Horizon 2020 Framework Programme for Research and Innovation. It is established by the Scientific Council of the ERC Grant Agreement no. 669792 (Allegro).
%\vspace*{\fill}

% in case there is an illustrator info on the illustrators
%%\vspace*{\fill}
%%\noindent \begin{minipage}[b]{1.0\columnwidth}%
%%\noindent Cover illustration: xxx\\
%%Cover design: xxx \end{minipage}%

%%% if this page should be empty, use the following command to
%%% circumvent de Acrobat error message about wrong contents object
%%% type:
%\vspace*{\fill}
%{\tiny .}


%%%%%%%%%%%%%%%%%%%%%%%%%%%%%%%%%%%%%%%%%%%%%%%%%%%%%%%%%%%%%%%%%%%%%
%% on page 3 the page as required by the "pedel"
%%%%%%%%%%%%%%%%%%%%%%%%%%%%%%%%%%%%%%%%%%%%%%%%%%%%%%%%%%%%%%%%%%%%%
\newpage
\thispagestyle{empty}
\vspace*{30mm}

\begin{center}
     \textbf{\huge \thetitle}\\[\baselineskip]\textbf{\LARGE \thesubtitle}
\end{center}

\vspace{20mm}

\begin{center}
    \textbf{Dissertation}
\end{center}
%\vspace{7mm}
\smallskip{}

\selectlanguage{american}

%\noindent{\Huge \textbf{Publications}}
\begin{small}

\begin{center}for the purpose of obtaining the degree of doctor\\
\vspace{1mm}
at Delft University of Technology\\
\vspace{1mm}
by the authority of the Rector Magnificus, Prof.dr.ir. T.H.J.J. van den Hagen\\
\vspace{1mm}
chair of the Board for Doctorates\\
\vspace{1mm}
to be defended publicly on\\
\vspace{1mm}
Friday 18th, September 2020 at 12:30 o'clock
\vspace{1mm}
\vspace{10mm} %5mm vertical space

by

\vspace{10mm}
\textbf{\theauthor}\\
\vspace{1mm}
Master of Science in Computer Science\\
\vspace{1mm}
Delft University of Technology\\
\vspace{1mm}
born in Jiajiang, Sichuan
\vspace{1mm}
\end{center}

\end{small}

% ----------------------------------------------
% Start, hidden title page for Graduate School

%\newpage
%\thispagestyle{empty}
%\vspace*{30mm}
%
%\begin{center}
%     \textbf{\huge \thetitle}\\[\baselineskip]\textbf{\LARGE \thesubtitle}
%\end{center}
%%\vspace{30mm}
%%\noindent \begin{center}\begin{tabbing}
%%\hspace*{70mm}\=\kill
%%\noindent Independent members:\\
%%Prof.dr.\ E.\ Simperl \> King's College London\\
%%Dr.\ M.\ Gonzalez \> University of California, Berkeley\\
%%Prof.dr.lr\ M.\ Van\ Steen \> University of Twente\\
%%Prof.dr.ir\ G.J.P.M.\ Houben \> Delft University of Technology\\
%%\end{tabbing}\end{center}
%
%\begin{center}
%    \textbf{Independent members}
%\end{center}\vspace{3mm}
%
%\begin{center}
%Prof.dr. E. Simperl, King's College London\\
%Dr. M. Gonzalez, University of California, Berkeley\\
%Prof.dr.lr M. Van Steen, University of Twente\\
%Prof.dr.ir G.J.P.M. Houben, Delft University of Technology\\
%\end{center}\vspace{3mm}
%
%\begin{center}
%    \textbf{Dissertation}
%\end{center}\vspace{3mm}
%
%%\smallskip{}
%
%\selectlanguage{american}
%
%%\noindent{\Huge \textbf{Publications}}
%\begin{small}
%
%\begin{center}for the purpose of obtaining the degree of doctor\\
%\vspace{1mm}
%at Delft University of Technology\\
%\vspace{1mm}
%by the authority of the Rector Magnificus, Prof.dr.ir. T.H.J.J. van den Hagen\\
%\vspace{1mm}
%chair of the Board for Doctorates\\
%\vspace{1mm}
%to be defended publicly on\\
%\vspace{1mm}
%Monday xx, Month 2020 at 00:00 o'clock
%\vspace{1mm}
%\vspace{4mm} %5mm vertical space
%
%by
%
%\vspace{4mm}
%\textbf{Vincent X. Gong}\\
%\vspace{1mm}
%Master of Science in Computer Science\\
%\vspace{1mm}
%Delft University of Technology\\
%\vspace{1mm}
%born in Jiajiang, Sichuan
%\vspace{1mm}
%\end{center}
%\end{small}

% end, hidden title page for Graduate School
% ----------------------------------------------

%%%%%%%%%%%%%%%%%%%%%%%%%%%%%%%%%%%%%%%%%%%%%%%%%%%%%%%%%%%%%%%%%%%%%%%%%%%%%%%%%%%%%%%%%
%% and on page 4 an overview of your review committee
%%%%%%%%%%%%%%%%%%%%%%%%%%%%%%%%%%%%%%%%%%%%%%%%%%%%%%%%%%%%%%%%%%%%%%%%%%%%%%%%%%%%%%%%%
\newpage %
\thispagestyle{empty}

\begin{small}

\noindent This dissertation has been approved by the promotors.\\
\bigskip{}

\noindent Composition of the doctoral committee:
%\smallskip{}
\noindent \begin{flushleft}\begin{tabbing}
\hspace*{70mm}\=\kill

Rector Magnificus \> Chairperson\\
% rem blow text until GS asks
Prof. dr. ir.\ S. P. Hoogendoorn \> Delft University of Technology, promotor\\
Dr. ir.\ W. Daamen \> Technische Universiteit Delft, promotor\\
Prof. dr. ir.\ A. Bozzon \> Technische Universiteit Delft, copromotor\\
% rem above text until GS asks
\bigskip{}\\

\noindent Independent members:\\
% rem blow text until GS asks
Prof. dr. ir.\ G. J. P. M.\ Houben \> Delft University of Technology\\
Prof. dr.\ E.\ Simperl \> King's College London, United Kingdom\\
Dr.\ M.\ Gonzalez \> University of California, Berkeley, United States\\
Prof. dr. ir.\ M.\ Van\ Steen \> University of Twente\\
% rem above text until GS asks
\end{tabbing}\end{flushleft}

\vspace*{\fill}
\selectlanguage{american}

% if you are publishing in the trail series, uncomment the following and include your number
%\noindent TRAIL Thesis Series T2007/14, The Netherlands TRAIL Research School

\bigskip{}
\noindent\textbf{TRAIL Thesis Series no. T2020/xx, the Netherlands Research School TRAIL}\\
\smallskip{}

\noindent TRAIL\\
P.O. BOX 5017\\
2600 GA Delft\\
The Netherlands\\
E-mail: info@rsTRAIL.nl\\
\bigskip{}

\noindent ISBN: your-isbn-number\\
\smallskip{}

%\noindent Copyright \copyright~2020 by \theauthor\\
\noindent Copyright \copyright~2020 by \theauthor\\
\smallskip{}

\noindent All rights reserved. No part of the material protected by this copyright notice may be reproduced or utilized in any form or by any means, electronic or mechanical, including photocopying, recording or by any information storage and retrieval system, without written permission of the author.

\smallskip{}
\noindent Printed in the Netherlands

\end{small}

%%%%%%%%%%%%%%%%%%%%%%%%%%%%%%%%%%%%%%%%%%%%%%%%%%%%%%%%%%%%%%%%%%%%
%% the back (page 5)
%%%%%%%%%%%%%%%%%%%%%%%%%%%%%%%%%%%%%%%%%%%%%%%%%%%%%%%%%%%%%%%%%%%%
%\newpage
%\thispagestyle{empty}
%
%\vspace*{\fill}
%\rightline{I will say something.}
%\vspace*{\fill}

% in case there is an illustrator info on the illustrators
%%\vspace*{\fill}
%%\noindent \begin{minipage}[b]{1.0\columnwidth}%
%%\noindent Cover illustration: xxx\\
%%Cover design: xxx \end{minipage}%

%%% if this page should be empty, use the following command to
%%% circumvent de Acrobat error message about wrong contents object
%%% type:
%\vspace*{\fill}
%{\tiny .}



%%% Local Variables: 
%%% mode: latex
%%% TeX-master: "../2007-thesis"
%%% End: 
}
% define relative path to images for this .tex file
\graphicspath{{./0-misc/images/}}


%%%%%%%%%%%%%%%%%%%%%%%%%%%%%%%%%%%%%%%%%%%%%%%%%%%%%%%%%%%%%%%%%%%%
% define below the title of your thesis, the subtitle, and your own name

\newcommand{\thetitle}{Using Social Media to Characterise Crowds in City Events for Crowd Management} 
\newcommand{\thesubtitle}{} 
\newcommand{\theauthor}{Vincent\ X.\ GONG}


%%%%%%%%%%%%%%%%%%%%%%%%%%%%%%%%%%%%%%%%%%%%%%%%%%%%%%%%%%%%%%%%%%%%
%% First the initial page with only the title and the author name
%%%%%%%%%%%%%%%%%%%%%%%%%%%%%%%%%%%%%%%%%%%%%%%%%%%%%%%%%%%%%%%%%%%%
\pagestyle{fancy} %% start page headers and footers
\thispagestyle{empty}
\vspace*{30mm}

\begin{center}
    \textbf{\huge \thetitle}\\[\baselineskip]\textbf{\LARGE \thesubtitle}
\end{center}


\vspace{40mm}

\begin{center}
    \Large\theauthor
\end{center}

%%%%%%%%%%%%%%%%%%%%%%%%%%%%%%%%%%%%%%%%%%%%%%%%%%%%%%%%%%%%%%%%%%%%
%% the back (page 2)
%%%%%%%%%%%%%%%%%%%%%%%%%%%%%%%%%%%%%%%%%%%%%%%%%%%%%%%%%%%%%%%%%%%%
\newpage

\thispagestyle{empty}

\vspace*{\fill}
% possibly include your logos and make acknowledgements to funding

\noindent \includegraphics[width=7em]{erc.png}\hspace*{0.75cm}
\includegraphics[width=7em]{allegro.png}\hspace*{0.75cm}
\includegraphics[width=12.5em]{ams.png}\hspace*{0.75cm}\\
%\noindent \includegraphics[width=21em]{logo_allegro}\\

\noindent The research leading to this dissertation has received funding from the European Research Council under the European Union Horizon 2020 Framework Programme for Research and Innovation. It is established by the Scientific Council of the ERC Grant Agreement no. 669792 (Allegro).
%\vspace*{\fill}

% in case there is an illustrator info on the illustrators
%%\vspace*{\fill}
%%\noindent \begin{minipage}[b]{1.0\columnwidth}%
%%\noindent Cover illustration: xxx\\
%%Cover design: xxx \end{minipage}%

%%% if this page should be empty, use the following command to
%%% circumvent de Acrobat error message about wrong contents object
%%% type:
%\vspace*{\fill}
%{\tiny .}


%%%%%%%%%%%%%%%%%%%%%%%%%%%%%%%%%%%%%%%%%%%%%%%%%%%%%%%%%%%%%%%%%%%%%
%% on page 3 the page as required by the "pedel"
%%%%%%%%%%%%%%%%%%%%%%%%%%%%%%%%%%%%%%%%%%%%%%%%%%%%%%%%%%%%%%%%%%%%%
\newpage
\thispagestyle{empty}
\vspace*{30mm}

\begin{center}
     \textbf{\huge \thetitle}\\[\baselineskip]\textbf{\LARGE \thesubtitle}
\end{center}

\vspace{20mm}

\begin{center}
    \textbf{Dissertation}
\end{center}
%\vspace{7mm}
\smallskip{}

\selectlanguage{american}

%\noindent{\Huge \textbf{Publications}}
\begin{small}

\begin{center}for the purpose of obtaining the degree of doctor\\
\vspace{1mm}
at Delft University of Technology\\
\vspace{1mm}
by the authority of the Rector Magnificus, Prof.dr.ir. T.H.J.J. van den Hagen\\
\vspace{1mm}
chair of the Board for Doctorates\\
\vspace{1mm}
to be defended publicly on\\
\vspace{1mm}
Friday 18th, September 2020 at 12:30 o'clock
\vspace{1mm}
\vspace{10mm} %5mm vertical space

by

\vspace{10mm}
\textbf{\theauthor}\\
\vspace{1mm}
Master of Science in Computer Science\\
\vspace{1mm}
Delft University of Technology\\
\vspace{1mm}
born in Jiajiang, Sichuan
\vspace{1mm}
\end{center}

\end{small}

% ----------------------------------------------
% Start, hidden title page for Graduate School

%\newpage
%\thispagestyle{empty}
%\vspace*{30mm}
%
%\begin{center}
%     \textbf{\huge \thetitle}\\[\baselineskip]\textbf{\LARGE \thesubtitle}
%\end{center}
%%\vspace{30mm}
%%\noindent \begin{center}\begin{tabbing}
%%\hspace*{70mm}\=\kill
%%\noindent Independent members:\\
%%Prof.dr.\ E.\ Simperl \> King's College London\\
%%Dr.\ M.\ Gonzalez \> University of California, Berkeley\\
%%Prof.dr.lr\ M.\ Van\ Steen \> University of Twente\\
%%Prof.dr.ir\ G.J.P.M.\ Houben \> Delft University of Technology\\
%%\end{tabbing}\end{center}
%
%\begin{center}
%    \textbf{Independent members}
%\end{center}\vspace{3mm}
%
%\begin{center}
%Prof.dr. E. Simperl, King's College London\\
%Dr. M. Gonzalez, University of California, Berkeley\\
%Prof.dr.lr M. Van Steen, University of Twente\\
%Prof.dr.ir G.J.P.M. Houben, Delft University of Technology\\
%\end{center}\vspace{3mm}
%
%\begin{center}
%    \textbf{Dissertation}
%\end{center}\vspace{3mm}
%
%%\smallskip{}
%
%\selectlanguage{american}
%
%%\noindent{\Huge \textbf{Publications}}
%\begin{small}
%
%\begin{center}for the purpose of obtaining the degree of doctor\\
%\vspace{1mm}
%at Delft University of Technology\\
%\vspace{1mm}
%by the authority of the Rector Magnificus, Prof.dr.ir. T.H.J.J. van den Hagen\\
%\vspace{1mm}
%chair of the Board for Doctorates\\
%\vspace{1mm}
%to be defended publicly on\\
%\vspace{1mm}
%Monday xx, Month 2020 at 00:00 o'clock
%\vspace{1mm}
%\vspace{4mm} %5mm vertical space
%
%by
%
%\vspace{4mm}
%\textbf{Vincent X. Gong}\\
%\vspace{1mm}
%Master of Science in Computer Science\\
%\vspace{1mm}
%Delft University of Technology\\
%\vspace{1mm}
%born in Jiajiang, Sichuan
%\vspace{1mm}
%\end{center}
%\end{small}

% end, hidden title page for Graduate School
% ----------------------------------------------

%%%%%%%%%%%%%%%%%%%%%%%%%%%%%%%%%%%%%%%%%%%%%%%%%%%%%%%%%%%%%%%%%%%%%%%%%%%%%%%%%%%%%%%%%
%% and on page 4 an overview of your review committee
%%%%%%%%%%%%%%%%%%%%%%%%%%%%%%%%%%%%%%%%%%%%%%%%%%%%%%%%%%%%%%%%%%%%%%%%%%%%%%%%%%%%%%%%%
\newpage %
\thispagestyle{empty}

\begin{small}

\noindent This dissertation has been approved by the promotors.\\
\bigskip{}

\noindent Composition of the doctoral committee:
%\smallskip{}
\noindent \begin{flushleft}\begin{tabbing}
\hspace*{70mm}\=\kill

Rector Magnificus \> Chairperson\\
% rem blow text until GS asks
Prof. dr. ir.\ S. P. Hoogendoorn \> Delft University of Technology, promotor\\
Dr. ir.\ W. Daamen \> Technische Universiteit Delft, promotor\\
Prof. dr. ir.\ A. Bozzon \> Technische Universiteit Delft, copromotor\\
% rem above text until GS asks
\bigskip{}\\

\noindent Independent members:\\
% rem blow text until GS asks
Prof. dr. ir.\ G. J. P. M.\ Houben \> Delft University of Technology\\
Prof. dr.\ E.\ Simperl \> King's College London, United Kingdom\\
Dr.\ M.\ Gonzalez \> University of California, Berkeley, United States\\
Prof. dr. ir.\ M.\ Van\ Steen \> University of Twente\\
% rem above text until GS asks
\end{tabbing}\end{flushleft}

\vspace*{\fill}
\selectlanguage{american}

% if you are publishing in the trail series, uncomment the following and include your number
%\noindent TRAIL Thesis Series T2007/14, The Netherlands TRAIL Research School

\bigskip{}
\noindent\textbf{TRAIL Thesis Series no. T2020/xx, the Netherlands Research School TRAIL}\\
\smallskip{}

\noindent TRAIL\\
P.O. BOX 5017\\
2600 GA Delft\\
The Netherlands\\
E-mail: info@rsTRAIL.nl\\
\bigskip{}

\noindent ISBN: your-isbn-number\\
\smallskip{}

%\noindent Copyright \copyright~2020 by \theauthor\\
\noindent Copyright \copyright~2020 by \theauthor\\
\smallskip{}

\noindent All rights reserved. No part of the material protected by this copyright notice may be reproduced or utilized in any form or by any means, electronic or mechanical, including photocopying, recording or by any information storage and retrieval system, without written permission of the author.

\smallskip{}
\noindent Printed in the Netherlands

\end{small}

%%%%%%%%%%%%%%%%%%%%%%%%%%%%%%%%%%%%%%%%%%%%%%%%%%%%%%%%%%%%%%%%%%%%
%% the back (page 5)
%%%%%%%%%%%%%%%%%%%%%%%%%%%%%%%%%%%%%%%%%%%%%%%%%%%%%%%%%%%%%%%%%%%%
%\newpage
%\thispagestyle{empty}
%
%\vspace*{\fill}
%\rightline{I will say something.}
%\vspace*{\fill}

% in case there is an illustrator info on the illustrators
%%\vspace*{\fill}
%%\noindent \begin{minipage}[b]{1.0\columnwidth}%
%%\noindent Cover illustration: xxx\\
%%Cover design: xxx \end{minipage}%

%%% if this page should be empty, use the following command to
%%% circumvent de Acrobat error message about wrong contents object
%%% type:
%\vspace*{\fill}
%{\tiny .}



%%% Local Variables: 
%%% mode: latex
%%% TeX-master: "../2007-thesis"
%%% End: 

\clearemptydoublepage

%%%%%%%%%%%%%%%%%%%%%%%%%%%%%%%%%%%%%%%%%%%%%%%%%%%%%%%%%%%%%%%%
\mainmatter

% here comes the preface

% to make sure that from here on american hyphenation is used, include
% the following

\selectlanguage{american}

\pagestyle{fancy}
\pagenumbering{roman}

% the 5 title pages defined in 0-misc/title.text have no numbers,
% so we manually set the next page, which should have a page number, to 6

\setcounter{page}{5}
%%%%%%%%%%%%%%%%%%%%%%%%%%%%%%%%%%%%%%%%%%%%%%%%%%%%%%%%%%%%%%%%

\clearemptydoublepage
%\dedication{\epigraph{I have something to say.}{}}
%\vspace*{-0.4cm}
%\vspace*{4cm}
\dedication{\epigraph{I have something to say.}{X. Gong}}

%\begin*{dedication}
%I have something to say.
%\end*{dedication}







\clearemptydoublepage

%%%%%%%%%%%%%%%%%%%%%%%%%%%%%%%%%%%%%%%%%%%%%%%%%%%%%%%%%%%%%%%%


\clearemptydoublepage

\chapter*{Acknowledgement}
% to make sure that TOC fits on 3 pages:
\addtocontents{toc}{\vspace*{-0.9cm}}
\addcontentsline{toc}{chapter}{Preface}

%\renewcommand{\imagepath}[1]{\fname{0-misc/images/#1}}
\graphicspath{{./0-misc/images/}}

\vspace*{-0.4cm}


Acknowledgement
\clearemptydoublepage

%%%%%%%%%%%%%%%%%%%%%%%%%%%%%%%%%%%%%%%%%%%%%%%%%%%%%%%%%%%%%%%%

% here comes the table of contents

\tableofcontents
\clearemptydoublepage

%%%%%%%%%%%%%%%%%%%%%%%%%%%%%%%%%%%%%%%%%%%%%%%%%%%%%%%%%%%%%%%%

%% reset numbering and set numbering to arabic
\pagenumbering{arabic}
\setcounter{page}{1}

%%%%%%%%%%%%%%%%%%%%%%%%%%%%%%%%%%%%%%%%%%%%%%%%%%%%%%%%%%%%%%%%

% here come the main chapters

% for reviewing purposes it may be useful to include line numbers.
% (this requires \usepackage{lineno} above to be uncommented)
\begin{linenumbers}[1]

% chapter 1
%%%%%%%%%%%%%%%%%%%%%%%%%%%%%%%%%%%%%%%%%%%%%%%%%%%%%%%%%%%%%%%%%%%%%%%

% make sure that "1-intro/images" refers to the subdirectory in which
% you store images for this chapter, relative to the base directory of
% your thesis (which contains the main thesis.tex file)

%\renewcommand{\imagepath}[1]{\fname{1-intro/images/#1}}
%\graphicspath{{./1-intro/images/}}

%%%%%%%%%%%%%%%%%%%%%%%%%%%%%%%%%%%%%%%%%%%%%%%%%%%%%%%%%%%%%%%%%%%%%%%

\chapter{Introduction}\label{ch1_title}

%In this chapter we present the background and the motivation for the
%research addressed in this thesis.  In Section \ref{transintro} we
%\ldots In Section \ref{overviewandroadmap} we conclude the chapter
%with an overview and road map of this thesis, and a list of the
%contributions to the state of the art.
%
%Parts of this chapter have been published in \cite{NegDeS:05-021}.

%\section{Background}\label{ch1_background}

City-scale events are getting more popular and attract a large number of people participating in various activities. For instance, on King’s Day, a national holiday in the Netherlands, a huge amount of people pour into the city and gather in the urban area, participating in various activities such as street parties, music festivals and boat parades. Event stakeholders, such as event organisers, police, municipalities, and crowd managers manage the crowd to avoid incidents. Crowd management practice consists of two phases \citep{martella2017current}, i.e. the planning phase and operational phase. In the planning phase, crowd managers require the past event data to infer guidelines, and to perform computer simulations of the crowds in the event.

\clearemptydoublepage 

% chapter 2
\chapter{Crowd Characterization for Crowd Management using Social Media Data in City Events}\label{ch2_title}

%\renewcommand{\imagepath}[1]{\fname{2-serpar/images/#1}}
%\graphicspath{{./2-serpar/images/}}
\noindent\rule{\textwidth}{1pt}
%\leftline{}

In this chapter, we characterise city events in terms of various aspects using social media data. This answers the first research question, i.e. \textbf{RQ1. To what extent social media data are able to characterize crowds in city events, in terms of demographic composition, city-role composition, spatio-temporal distribution, Points of Interest preferences and word use?} 

To this end, we screen a set of factors (i.e. visitor profile, crowd size, density, mobility, location, and semantics) that characterize crowd behaviour and introduce a set of proxies (i.e. demographics, city-role, crowd temporal distribution, post position, Points of Interests, and word use) derived from social media data. Furthermore, we characterize the crowd in two city-scale events, Sail 2015 and King's Day 2016, in terms of these proxies, and comparing them with information collected from events organizers and programs.

Our findings show that it is possible to characterize crowds in city-scale events using social media data, thus paving the way for new real-time and planning applications on crowd monitoring and management for city-scale events.


%This chapter is currently under review for journal publication.
This chapter is published as a journal article: Gong, V. X., Daamen, W., Bozzon, A., \& Hoogendoorn, S. P. (2020). Crowd characterization for crowd management using social media data in city events. \emph{Travel Behaviour and Society}, 20, 192-212.

\noindent\rule{\textwidth}{1pt}

\newpage


\vspace*{-10mm}
\section{Introduction}\label{ch2_2.1_introduction}

As cities compete for global importance and influence, city-scale public events are becoming an important ingredient to foster tourism and economic growth. Sports events, thematic exhibitions, and national celebrations are examples of city-scale events that take place in vast urban areas, and attract large amounts of participants within short time spans. The scale and intensity of these happenings demand technological solutions supporting stakeholders (e.g. event organizers, public and safety authorities, attendees) to monitor and manage the crowd.

\clearemptydoublepage

\chapter{Conclusions, implications and recommendations}\label{ch6_title}

In this chapter, we present our main findings and conclusions for each research question, followed by the overall conclusions and implications for practice. Finally, we provide recommendations for future research.


%\clearemptydoublepage


%%%%%%%%%%%%%%%%%%%%%%%%%%%%%%%%%%%%%%%%%%%%%%%%%%%%%%%%%%%%%%%%

% here come the appendices, if any

%\appendix

%\chapter{}\label{tabs_and_cal}

%\renewcommand{\imagepath}[1]{\fname{A1-my_appendix/images/#1}}
%\graphicspath{{./A1-my_appendix/images/}}

%Here comes the appendix.

% LocalWords:  subnetworks subnetwork


\section{Social media geo-posts sent around Sail and Kingsday event}

\begin{itemize}
\item The temporal distribution of posts sent by people observed from social media in Figure \ref{fig:temperal_distribution_app}.
\item The Total, Max, and Minimal number of geo-posts sent around Sail event 2015 in Table \ref{tab:temporal_sail_app} and around King's Day event 2016 in Table \ref{tab:temporal_kd_app}.
% \item The Total, Max, and Minimal Number of Geo-posts Sent Around King's Day Event 2016 in Table \ref{tab:temporal_kd_app}.
\end{itemize}


\begin{figure*}
\centering     %%% not \center
% \resizebox{\textwidth}{!}{%
\subfigure[Sail]{\label{fig:sail_long}\includegraphics[width=0.88\linewidth]{./2-chara/images/ch3/time/sail_long}}
\subfigure[Sail, Twitter]{\label{fig:sail_long_twitter}\includegraphics[width=0.88\linewidth]{./2-chara/images/ch3/time/sail_long_twitter}}

\subfigure[Kingsday]{\label{fig:kd_long}\includegraphics[width=0.88\linewidth]{./2-chara/images/ch3/time/kd_long}}
\subfigure[Kingsday, Twitter]{\label{fig:kd_long_twitter}\includegraphics[width=0.88\linewidth]{./2-chara/images/ch3/time/kd_long_twitter}}
% }
\caption{The temporal distribution of posts sent by people observed from social media}
\label{fig:temperal_distribution_app}
\end{figure*}


% start: two tables for the long diagram

\begin{landscape}
\begin{table}[]
% \centering
\caption{The Total, Max, and Minimal Number of Geo-posts Sent Around Sail Event 2015}
\begin{center}
\label{tab:temporal_sail_app}
\begin{threeparttable}
\resizebox{1.37\textwidth}{!}{%
\begin{tabular}{llrrrrrrlrlrrrrrlrlrrlrrrrr} 
\toprule
\multicolumn{2}{l}{\multirow{5}{*}{}} & \multicolumn{25}{c}{Date} \\
\multicolumn{2}{l}{} & \multicolumn{25}{l}{} \\
\multicolumn{2}{l}{} & \multicolumn{1}{c}{} & \multicolumn{1}{c}{} & \multicolumn{1}{c}{} & \multicolumn{1}{c}{} & \multicolumn{1}{c}{} & \multicolumn{1}{c}{} & \multicolumn{1}{r}{} & \multicolumn{1}{l}{Weekend} & \multicolumn{1}{r}{} & \multicolumn{5}{c}{Sail Event} &  & \multicolumn{1}{l}{} & \multicolumn{1}{r}{} & \multicolumn{2}{c}{Weekend} &  & \multicolumn{1}{c}{} & \multicolumn{1}{c}{} & \multicolumn{1}{c}{} & \multicolumn{1}{c}{} & \multicolumn{1}{c}{} \\ 
\cmidrule{10-10}\cmidrule{12-16}\cmidrule{20-21}
\multicolumn{2}{l}{} & \multicolumn{1}{c}{Sun} & \multicolumn{1}{c}{Mon} & \multicolumn{1}{c}{Tue} & \multicolumn{1}{c}{Wed} & \multicolumn{1}{c}{Thu} & \multicolumn{1}{c}{Fri} & \multicolumn{1}{r}{} & \multicolumn{1}{c}{Sat} & \multicolumn{1}{r}{} & \multicolumn{1}{c}{Sun} & \multicolumn{1}{c}{Mon} & \multicolumn{1}{c}{Tue} & \multicolumn{1}{c}{Wed} & \multicolumn{1}{c}{Thu} & \multicolumn{1}{r}{} & \multicolumn{1}{c}{Fri} & \multicolumn{1}{r}{} & \multicolumn{1}{c}{Sat} & \multicolumn{1}{c}{Sun} & \multicolumn{1}{r}{} & \multicolumn{1}{c}{Mon} & \multicolumn{1}{c}{Tue} & \multicolumn{1}{c}{Wed} & \multicolumn{1}{c}{Thu} & \multicolumn{1}{c}{Fri} \\
\multicolumn{2}{l}{} & \multicolumn{1}{l}{Aug 12} & Aug 13 & Aug 14 & Aug 15 & Aug 16 & Aug 17 & \multicolumn{1}{r}{} & \multicolumn{1}{l}{Aug 18} & \multicolumn{1}{r}{} & \multicolumn{1}{l}{Aug 19} & Aug 20 & Aug 21 & Aug 22 & Aug 23 & \multicolumn{1}{r}{} & \multicolumn{1}{l}{Aug 24} & \multicolumn{1}{r}{} & \multicolumn{1}{l}{Aug 25} & Aug 26 & \multicolumn{1}{r}{} & \multicolumn{1}{l}{Aug 27} & Aug 28 & Aug 29 & Aug 30 & Aug 31 \\ 
\midrule
\multirow{5}{*}{Twitter} & \#P. of day & 759 & 686 & 727 & 757 & 591 & 629 &  & 754 &  & 1083 & 1092 & 1154 & 1212 & 1002 &  & 892 &  & 834 & 868 &  & 832 & 926 & 905 & 980 & 836 \\
 & Max \#P., Time (hour) & 15:00 & 18:00 & 20:00 & 19:00 & 15:00 & 19:00 &  & 15:00 &  & 14:00 & 18:00 & 20:00 & 14:00 & 14:00 &  & 19:00 &  & 16:00 & 17:00 &  & 13:00 & 20:00 & 12:00 & 13:00 & 17:00 \\
 & Max \#P. & 55 & 51 & 63 & 79 & 54 & 55 &  & 61 &  & 101 & 84 & 102 & 109 & 90 &  & 79 &  & 60 & 66 &  & 75 & 83 & 74 & 78 & 63 \\
 & Min \#P., Time (hour) & 04:00 & 05:00 & 05:00 & 06:00 & 04:00 & 05:00 &  & 03:00 &  & 03:00 & 04:00 & 05:00 & 05:00 & 05:00 &  & 05:00 &  & 03:00 & 03:00 &  & 05:00 & 04:00 & 04:00 & 06:00 & 05:00 \\
 & Min \#P. & 3 & 0 & 1 & 3 & 0 & 1 &  & 2 &  & 2 & 1 & 4 & 4 & 3 &  & 3 &  & 0 & 2 &  & 0 & 4 & 2 & 3 & 2 \\ 
\midrule
\multirow{5}{*}{Instagram} & \#P. of day & 1291 & 1435 & 1510 & 1696 & 1943 & 2012 &  & 2796 &  & 3775 & 2967 & 2766 & 3429 & 4638 &  & 3278 &  & 1796 & 1751 &  & 1815 & 2156 & 2402 & 2754 & 2268 \\
 & Max \#P., Time (hour) & 19:00 & 21:00 & 21:00 & 20:00 & 19:00 & 22:00 &  & 20:00 &  & 20:00 & 21:00 & 20:00 & 20:00 & 20:00 &  & 18:00 &  & 20:00 & 18:00 &  & 22:00 & 19:00 & 19:00 & 18:00 & 17:00 \\
 & Max \#P. & 112 & 121 & 119 & 139 & 143 & 166 &  & 261 &  & 355 & 210 & 235 & 299 & 355 &  & 264 &  & 132 & 135 &  & 135 & 156 & 172 & 201 & 157 \\
 & Min \#P., Time (hour) & 03:00 & 05:00 & 05:00 & 05:00 & 05:00 & 04:00 &  & 05:00 &  & 04:00 & 04:00 & 06:00 & 06:00 & 05:00 &  & 05:00 &  & 04:00 & 06:00 &  & 05:00 & 06:00 & 05:00 & 06:00 & 05:00 \\
 & Min \#P., Time (hour) & 5 & 4 & 4 & 5 & 3 & 5 &  & 5 &  & 6 & 12 & 9 & 8 & 11 &  & 16 &  & 9 & 10 &  & 10 & 16 & 6 & 12 & 9 \\
\bottomrule
\end{tabular}
}

\begin{tablenotes}
      \scriptsize
      \item \textbf{\#GP.}: amount of Geo-posts.
      \item \textbf{Max \#GP. Time(hour)}: the time in hour during which the max amount of Geo-posts is observed.
\end{tablenotes}
\end{threeparttable}
\end{center}
\end{table}
% \end{landscape}



% \begin{landscape}
\begin{table}[]
% \centering
\caption{The Total, Max, and Minimal Number of Geo-posts Sent Around King's Day Event 2016}
\begin{center}
\label{tab:temporal_kd_app}
\begin{threeparttable}
\resizebox{1.37\textwidth}{!}{%
\begin{tabular}{llrrrlrrlrlrrlrrlrrlrrrr} 
\toprule
\multicolumn{2}{l}{\multirow{5}{*}{}} & \multicolumn{22}{c}{Date} \\
\multicolumn{2}{l}{} & \multicolumn{22}{l}{} \\
\multicolumn{2}{l}{} & \multicolumn{1}{c}{} & \multicolumn{1}{c}{} & \multicolumn{1}{c}{} &  & \multicolumn{2}{c}{Weekend} &  & \multicolumn{1}{l}{} &  & \multicolumn{1}{l}{K. Night} & \multicolumn{1}{l}{K. Day} &  & \multicolumn{1}{l}{} & \multicolumn{1}{l}{} &  & \multicolumn{2}{c}{Weekend} &  & \multicolumn{1}{c}{} & \multicolumn{1}{c}{} & \multicolumn{1}{c}{} & \multicolumn{1}{c}{} \\ 
\cmidrule{7-8}\cmidrule{12-13}\cmidrule{18-19}
\multicolumn{2}{l}{} & \multicolumn{1}{c}{Wed} & \multicolumn{1}{c}{Thu} & \multicolumn{1}{c}{Fri} &  & \multicolumn{1}{c}{Sat} & \multicolumn{1}{c}{Sun} &  & \multicolumn{1}{c}{Mon} &  & \multicolumn{1}{c}{Tue} & \multicolumn{1}{c}{Wed} &  & \multicolumn{1}{c}{Thu} & \multicolumn{1}{c}{Fri} &  & \multicolumn{1}{c}{Sat} & \multicolumn{1}{c}{Sun} &  & \multicolumn{1}{c}{Mon} & \multicolumn{1}{c}{Tue} & \multicolumn{1}{c}{Wed} & \multicolumn{1}{c}{Thu} \\
\multicolumn{2}{l}{} & \multicolumn{1}{l}{April 20} & April 21 & April 22 &  & April 23 & April 24 &  & April 25 &  & April 26 & April 27 &  & April 28 & April 29 &  & April 30 & May 1 &  & May 2 & May 3 & May 4 & May 5 \\ 
\midrule
\multirow{5}{*}{Twitter} & \#P. of day & 822 & 926 & 839 &  & 918 & 886 &  & 768 &  & 855 & 1134 &  & 914 & 833 &  & 814 & 957 &  & 785 & 715 & 794 & 909 \\
 & Max \#P., Time (hour) & 16:00 & 22:00 & 16:00 &  & 20:00 & 13:00 &  & 17:00 &  & 23:00 & 16:00 &  & 16:00 & 16:00 &  & 14:00 & 16:00 &  & 17:00 & 18:00 & 18:00 & 14:00 \\
 & Max \#P. & 62 & 84 & 73 &  & 79 & 81 &  & 61 &  & 72 & 99 &  & 67 & 67 &  & 64 & 75 &  & 63 & 60 & 63 & 88 \\
 & Min \#P., Time (hour) & 01:00 & 04:00 & 05:00 &  & 05:00 & 04:00 &  & 05:00 &  & 05:00 & 06:00 &  & 04:00 & 06:00 &  & 05:00 & 04:00 &  & 04:00 & 03:00 & 04:00 & 05:00 \\
 & Min \#P. & 0 & 2 & 2 &  & 2 & 1 &  & 3 &  & 0 & 9 &  & 2 & 10 &  & 1 & 0 &  & 2 & 0 & 6 & 0 \\ 
\midrule
\multirow{5}{*}{Instagram} & \#P. of day & 3369 & 3402 & 3325 &  & 4006 & 4574 &  & 4277 &  & 5450 & 8902 &  & 5575 & 3407 &  & 3622 & 4160 &  & 3666 & 3647 & 3723 & 4959 \\
 & Max \#P., Time (hour) & 19:00 & 19:00 & 19:00 &  & 18:00 & 19:00 &  & 21:00 &  & 20:00 & 22:00 &  & 13:00 & 18:00 &  & 18:00 & 19:00 &  & 21:00 & 20:00 & 22:00 & 20:00 \\
 & Max \#P. & 248 & 256 & 238 &  & 299 & 326 &  & 339 &  & 425 & 715 &  & 394 & 225 &  & 271 & 355 &  & 276 & 289 & 276 & 396 \\
 & Min \#P., Time (hour) & 03:00 & 04:00 & 05:00 &  & 05:00 & 05:00 &  & 04:00 &  & 05:00 & 05:00 &  & 05:00 & 06:00 &  & 05:00 & 05:00 &  & 05:00 & 05:00 & 05:00 & 04:00 \\
 & Min \#P. & 19 & 19 & 14 &  & 14 & 22 &  & 17 &  & 14 & 48 &  & 36 & 20 &  & 15 & 16 &  & 20 & 18 & 20 & 18 \\
\bottomrule
\end{tabular}
}
\end{threeparttable}
\end{center}
\end{table}
\end{landscape}


\section{Areas of Interest based on clustering social media posts}

The Area of Interest detected by clustering PoIs using social media posts for 4 sub-events listed in Table 8 are shown in Figure \ref{fig:apd_poi_sail_parade}, \ref{fig:apd_poi_sail_fireworks}, \ref{fig:apd_poi_kingsday_parade}, and \ref{fig:apd_poi_kingsday_night}. These cluster maps are generated using DBSCAN algorithm \citep{ester1996density}, which groups points within distance of 50 meters (epsilon=0.05). Such groups will be identified as a cluster if the number of points in a group is more than 15 (min\_Points = 15), otherwise each points will be determined as outliers. Each colour represents one cluster. Black points are outliers which fail to form group with any other points. The count number denotes the amount of points in each cluster.

\clearpage

\begin{figure*}[]
\centering
% \captionsetup{justification=centering}
\includegraphics[width=0.8\linewidth]{./2-chara/images/ch3/poi/raw/sail_parade_clusters}
\caption{The Area of Interest detected by clustering PoIs using social media posts in Sail, Sail-in Parade}
\label{fig:apd_poi_sail_parade}
\end{figure*}

\begin{figure*}[]
\centering
% \captionsetup{justification=centering}
\includegraphics[width=0.8\linewidth]{./2-chara/images/ch3/poi/raw/sail_fireworks_clusters}
\caption{The Area of Interest detected by clustering PoIs using social media posts in Sail, Fireworks}
\label{fig:apd_poi_sail_fireworks}
\end{figure*}

\begin{figure*}[]
\centering
% \captionsetup{justification=centering}
\includegraphics[width=0.8\linewidth]{./2-chara/images/ch3/poi/raw/kingsday_parade_clusters}
\caption{The Area of Interest detected by clustering PoIs using social media posts in King's Day, Boat Parade}
\label{fig:apd_poi_kingsday_parade}
\end{figure*}

\begin{figure*}[]
\centering
% \captionsetup{justification=centering}
\includegraphics[width=0.8\linewidth]{./2-chara/images/ch3/poi/raw/kingsday_kingsnight_clusters}
\caption{The Area of Interest detected by clustering PoIs using social media posts in King's Day, King's Night}
\label{fig:apd_poi_kingsday_night}
\end{figure*}}
%\chapter{}\label{tabs_and_cal}

%\renewcommand{\imagepath}[1]{\fname{A1-my_appendix/images/#1}}
%\graphicspath{{./A1-my_appendix/images/}}

%Here comes the appendix.

% LocalWords:  subnetworks subnetwork


\section{Social media geo-posts sent around Sail and Kingsday event}

\begin{itemize}
\item The temporal distribution of posts sent by people observed from social media in Figure \ref{fig:temperal_distribution_app}.
\item The Total, Max, and Minimal number of geo-posts sent around Sail event 2015 in Table \ref{tab:temporal_sail_app} and around King's Day event 2016 in Table \ref{tab:temporal_kd_app}.
% \item The Total, Max, and Minimal Number of Geo-posts Sent Around King's Day Event 2016 in Table \ref{tab:temporal_kd_app}.
\end{itemize}


\begin{figure*}
\centering     %%% not \center
% \resizebox{\textwidth}{!}{%
\subfigure[Sail]{\label{fig:sail_long}\includegraphics[width=0.88\linewidth]{./2-chara/images/ch3/time/sail_long}}
\subfigure[Sail, Twitter]{\label{fig:sail_long_twitter}\includegraphics[width=0.88\linewidth]{./2-chara/images/ch3/time/sail_long_twitter}}

\subfigure[Kingsday]{\label{fig:kd_long}\includegraphics[width=0.88\linewidth]{./2-chara/images/ch3/time/kd_long}}
\subfigure[Kingsday, Twitter]{\label{fig:kd_long_twitter}\includegraphics[width=0.88\linewidth]{./2-chara/images/ch3/time/kd_long_twitter}}
% }
\caption{The temporal distribution of posts sent by people observed from social media}
\label{fig:temperal_distribution_app}
\end{figure*}


% start: two tables for the long diagram

\begin{landscape}
\begin{table}[]
% \centering
\caption{The Total, Max, and Minimal Number of Geo-posts Sent Around Sail Event 2015}
\begin{center}
\label{tab:temporal_sail_app}
\begin{threeparttable}
\resizebox{1.37\textwidth}{!}{%
\begin{tabular}{llrrrrrrlrlrrrrrlrlrrlrrrrr} 
\toprule
\multicolumn{2}{l}{\multirow{5}{*}{}} & \multicolumn{25}{c}{Date} \\
\multicolumn{2}{l}{} & \multicolumn{25}{l}{} \\
\multicolumn{2}{l}{} & \multicolumn{1}{c}{} & \multicolumn{1}{c}{} & \multicolumn{1}{c}{} & \multicolumn{1}{c}{} & \multicolumn{1}{c}{} & \multicolumn{1}{c}{} & \multicolumn{1}{r}{} & \multicolumn{1}{l}{Weekend} & \multicolumn{1}{r}{} & \multicolumn{5}{c}{Sail Event} &  & \multicolumn{1}{l}{} & \multicolumn{1}{r}{} & \multicolumn{2}{c}{Weekend} &  & \multicolumn{1}{c}{} & \multicolumn{1}{c}{} & \multicolumn{1}{c}{} & \multicolumn{1}{c}{} & \multicolumn{1}{c}{} \\ 
\cmidrule{10-10}\cmidrule{12-16}\cmidrule{20-21}
\multicolumn{2}{l}{} & \multicolumn{1}{c}{Sun} & \multicolumn{1}{c}{Mon} & \multicolumn{1}{c}{Tue} & \multicolumn{1}{c}{Wed} & \multicolumn{1}{c}{Thu} & \multicolumn{1}{c}{Fri} & \multicolumn{1}{r}{} & \multicolumn{1}{c}{Sat} & \multicolumn{1}{r}{} & \multicolumn{1}{c}{Sun} & \multicolumn{1}{c}{Mon} & \multicolumn{1}{c}{Tue} & \multicolumn{1}{c}{Wed} & \multicolumn{1}{c}{Thu} & \multicolumn{1}{r}{} & \multicolumn{1}{c}{Fri} & \multicolumn{1}{r}{} & \multicolumn{1}{c}{Sat} & \multicolumn{1}{c}{Sun} & \multicolumn{1}{r}{} & \multicolumn{1}{c}{Mon} & \multicolumn{1}{c}{Tue} & \multicolumn{1}{c}{Wed} & \multicolumn{1}{c}{Thu} & \multicolumn{1}{c}{Fri} \\
\multicolumn{2}{l}{} & \multicolumn{1}{l}{Aug 12} & Aug 13 & Aug 14 & Aug 15 & Aug 16 & Aug 17 & \multicolumn{1}{r}{} & \multicolumn{1}{l}{Aug 18} & \multicolumn{1}{r}{} & \multicolumn{1}{l}{Aug 19} & Aug 20 & Aug 21 & Aug 22 & Aug 23 & \multicolumn{1}{r}{} & \multicolumn{1}{l}{Aug 24} & \multicolumn{1}{r}{} & \multicolumn{1}{l}{Aug 25} & Aug 26 & \multicolumn{1}{r}{} & \multicolumn{1}{l}{Aug 27} & Aug 28 & Aug 29 & Aug 30 & Aug 31 \\ 
\midrule
\multirow{5}{*}{Twitter} & \#P. of day & 759 & 686 & 727 & 757 & 591 & 629 &  & 754 &  & 1083 & 1092 & 1154 & 1212 & 1002 &  & 892 &  & 834 & 868 &  & 832 & 926 & 905 & 980 & 836 \\
 & Max \#P., Time (hour) & 15:00 & 18:00 & 20:00 & 19:00 & 15:00 & 19:00 &  & 15:00 &  & 14:00 & 18:00 & 20:00 & 14:00 & 14:00 &  & 19:00 &  & 16:00 & 17:00 &  & 13:00 & 20:00 & 12:00 & 13:00 & 17:00 \\
 & Max \#P. & 55 & 51 & 63 & 79 & 54 & 55 &  & 61 &  & 101 & 84 & 102 & 109 & 90 &  & 79 &  & 60 & 66 &  & 75 & 83 & 74 & 78 & 63 \\
 & Min \#P., Time (hour) & 04:00 & 05:00 & 05:00 & 06:00 & 04:00 & 05:00 &  & 03:00 &  & 03:00 & 04:00 & 05:00 & 05:00 & 05:00 &  & 05:00 &  & 03:00 & 03:00 &  & 05:00 & 04:00 & 04:00 & 06:00 & 05:00 \\
 & Min \#P. & 3 & 0 & 1 & 3 & 0 & 1 &  & 2 &  & 2 & 1 & 4 & 4 & 3 &  & 3 &  & 0 & 2 &  & 0 & 4 & 2 & 3 & 2 \\ 
\midrule
\multirow{5}{*}{Instagram} & \#P. of day & 1291 & 1435 & 1510 & 1696 & 1943 & 2012 &  & 2796 &  & 3775 & 2967 & 2766 & 3429 & 4638 &  & 3278 &  & 1796 & 1751 &  & 1815 & 2156 & 2402 & 2754 & 2268 \\
 & Max \#P., Time (hour) & 19:00 & 21:00 & 21:00 & 20:00 & 19:00 & 22:00 &  & 20:00 &  & 20:00 & 21:00 & 20:00 & 20:00 & 20:00 &  & 18:00 &  & 20:00 & 18:00 &  & 22:00 & 19:00 & 19:00 & 18:00 & 17:00 \\
 & Max \#P. & 112 & 121 & 119 & 139 & 143 & 166 &  & 261 &  & 355 & 210 & 235 & 299 & 355 &  & 264 &  & 132 & 135 &  & 135 & 156 & 172 & 201 & 157 \\
 & Min \#P., Time (hour) & 03:00 & 05:00 & 05:00 & 05:00 & 05:00 & 04:00 &  & 05:00 &  & 04:00 & 04:00 & 06:00 & 06:00 & 05:00 &  & 05:00 &  & 04:00 & 06:00 &  & 05:00 & 06:00 & 05:00 & 06:00 & 05:00 \\
 & Min \#P., Time (hour) & 5 & 4 & 4 & 5 & 3 & 5 &  & 5 &  & 6 & 12 & 9 & 8 & 11 &  & 16 &  & 9 & 10 &  & 10 & 16 & 6 & 12 & 9 \\
\bottomrule
\end{tabular}
}

\begin{tablenotes}
      \scriptsize
      \item \textbf{\#GP.}: amount of Geo-posts.
      \item \textbf{Max \#GP. Time(hour)}: the time in hour during which the max amount of Geo-posts is observed.
\end{tablenotes}
\end{threeparttable}
\end{center}
\end{table}
% \end{landscape}



% \begin{landscape}
\begin{table}[]
% \centering
\caption{The Total, Max, and Minimal Number of Geo-posts Sent Around King's Day Event 2016}
\begin{center}
\label{tab:temporal_kd_app}
\begin{threeparttable}
\resizebox{1.37\textwidth}{!}{%
\begin{tabular}{llrrrlrrlrlrrlrrlrrlrrrr} 
\toprule
\multicolumn{2}{l}{\multirow{5}{*}{}} & \multicolumn{22}{c}{Date} \\
\multicolumn{2}{l}{} & \multicolumn{22}{l}{} \\
\multicolumn{2}{l}{} & \multicolumn{1}{c}{} & \multicolumn{1}{c}{} & \multicolumn{1}{c}{} &  & \multicolumn{2}{c}{Weekend} &  & \multicolumn{1}{l}{} &  & \multicolumn{1}{l}{K. Night} & \multicolumn{1}{l}{K. Day} &  & \multicolumn{1}{l}{} & \multicolumn{1}{l}{} &  & \multicolumn{2}{c}{Weekend} &  & \multicolumn{1}{c}{} & \multicolumn{1}{c}{} & \multicolumn{1}{c}{} & \multicolumn{1}{c}{} \\ 
\cmidrule{7-8}\cmidrule{12-13}\cmidrule{18-19}
\multicolumn{2}{l}{} & \multicolumn{1}{c}{Wed} & \multicolumn{1}{c}{Thu} & \multicolumn{1}{c}{Fri} &  & \multicolumn{1}{c}{Sat} & \multicolumn{1}{c}{Sun} &  & \multicolumn{1}{c}{Mon} &  & \multicolumn{1}{c}{Tue} & \multicolumn{1}{c}{Wed} &  & \multicolumn{1}{c}{Thu} & \multicolumn{1}{c}{Fri} &  & \multicolumn{1}{c}{Sat} & \multicolumn{1}{c}{Sun} &  & \multicolumn{1}{c}{Mon} & \multicolumn{1}{c}{Tue} & \multicolumn{1}{c}{Wed} & \multicolumn{1}{c}{Thu} \\
\multicolumn{2}{l}{} & \multicolumn{1}{l}{April 20} & April 21 & April 22 &  & April 23 & April 24 &  & April 25 &  & April 26 & April 27 &  & April 28 & April 29 &  & April 30 & May 1 &  & May 2 & May 3 & May 4 & May 5 \\ 
\midrule
\multirow{5}{*}{Twitter} & \#P. of day & 822 & 926 & 839 &  & 918 & 886 &  & 768 &  & 855 & 1134 &  & 914 & 833 &  & 814 & 957 &  & 785 & 715 & 794 & 909 \\
 & Max \#P., Time (hour) & 16:00 & 22:00 & 16:00 &  & 20:00 & 13:00 &  & 17:00 &  & 23:00 & 16:00 &  & 16:00 & 16:00 &  & 14:00 & 16:00 &  & 17:00 & 18:00 & 18:00 & 14:00 \\
 & Max \#P. & 62 & 84 & 73 &  & 79 & 81 &  & 61 &  & 72 & 99 &  & 67 & 67 &  & 64 & 75 &  & 63 & 60 & 63 & 88 \\
 & Min \#P., Time (hour) & 01:00 & 04:00 & 05:00 &  & 05:00 & 04:00 &  & 05:00 &  & 05:00 & 06:00 &  & 04:00 & 06:00 &  & 05:00 & 04:00 &  & 04:00 & 03:00 & 04:00 & 05:00 \\
 & Min \#P. & 0 & 2 & 2 &  & 2 & 1 &  & 3 &  & 0 & 9 &  & 2 & 10 &  & 1 & 0 &  & 2 & 0 & 6 & 0 \\ 
\midrule
\multirow{5}{*}{Instagram} & \#P. of day & 3369 & 3402 & 3325 &  & 4006 & 4574 &  & 4277 &  & 5450 & 8902 &  & 5575 & 3407 &  & 3622 & 4160 &  & 3666 & 3647 & 3723 & 4959 \\
 & Max \#P., Time (hour) & 19:00 & 19:00 & 19:00 &  & 18:00 & 19:00 &  & 21:00 &  & 20:00 & 22:00 &  & 13:00 & 18:00 &  & 18:00 & 19:00 &  & 21:00 & 20:00 & 22:00 & 20:00 \\
 & Max \#P. & 248 & 256 & 238 &  & 299 & 326 &  & 339 &  & 425 & 715 &  & 394 & 225 &  & 271 & 355 &  & 276 & 289 & 276 & 396 \\
 & Min \#P., Time (hour) & 03:00 & 04:00 & 05:00 &  & 05:00 & 05:00 &  & 04:00 &  & 05:00 & 05:00 &  & 05:00 & 06:00 &  & 05:00 & 05:00 &  & 05:00 & 05:00 & 05:00 & 04:00 \\
 & Min \#P. & 19 & 19 & 14 &  & 14 & 22 &  & 17 &  & 14 & 48 &  & 36 & 20 &  & 15 & 16 &  & 20 & 18 & 20 & 18 \\
\bottomrule
\end{tabular}
}
\end{threeparttable}
\end{center}
\end{table}
\end{landscape}


\section{Areas of Interest based on clustering social media posts}

The Area of Interest detected by clustering PoIs using social media posts for 4 sub-events listed in Table 8 are shown in Figure \ref{fig:apd_poi_sail_parade}, \ref{fig:apd_poi_sail_fireworks}, \ref{fig:apd_poi_kingsday_parade}, and \ref{fig:apd_poi_kingsday_night}. These cluster maps are generated using DBSCAN algorithm \citep{ester1996density}, which groups points within distance of 50 meters (epsilon=0.05). Such groups will be identified as a cluster if the number of points in a group is more than 15 (min\_Points = 15), otherwise each points will be determined as outliers. Each colour represents one cluster. Black points are outliers which fail to form group with any other points. The count number denotes the amount of points in each cluster.

\clearpage

\begin{figure*}[]
\centering
% \captionsetup{justification=centering}
\includegraphics[width=0.8\linewidth]{./2-chara/images/ch3/poi/raw/sail_parade_clusters}
\caption{The Area of Interest detected by clustering PoIs using social media posts in Sail, Sail-in Parade}
\label{fig:apd_poi_sail_parade}
\end{figure*}

\begin{figure*}[]
\centering
% \captionsetup{justification=centering}
\includegraphics[width=0.8\linewidth]{./2-chara/images/ch3/poi/raw/sail_fireworks_clusters}
\caption{The Area of Interest detected by clustering PoIs using social media posts in Sail, Fireworks}
\label{fig:apd_poi_sail_fireworks}
\end{figure*}

\begin{figure*}[]
\centering
% \captionsetup{justification=centering}
\includegraphics[width=0.8\linewidth]{./2-chara/images/ch3/poi/raw/kingsday_parade_clusters}
\caption{The Area of Interest detected by clustering PoIs using social media posts in King's Day, Boat Parade}
\label{fig:apd_poi_kingsday_parade}
\end{figure*}

\begin{figure*}[]
\centering
% \captionsetup{justification=centering}
\includegraphics[width=0.8\linewidth]{./2-chara/images/ch3/poi/raw/kingsday_kingsnight_clusters}
\caption{The Area of Interest detected by clustering PoIs using social media posts in King's Day, King's Night}
\label{fig:apd_poi_kingsday_night}
\end{figure*}
%\clearemptydoublepage


%%%%%%%%%%%%%%%%%%%%%%%%%%%%%%%%%%%%%%%%%%%%%%%%%%%%%%%%%%%%%%%%

% here comes the bibliography

%% next line is necessary to make link in table of contents appear correctly
\phantomsection
\addcontentsline{toc}{chapter}{Bibliography}

%\bibliographystyle{plainnat}
\bibliographystyle{TRAILbib}


% add here your own bibliography files
\bibliography{./references}
}
% next line is necessary to make link in table of contents appear correctly
\phantomsection
\addcontentsline{toc}{chapter}{Bibliography}

%\bibliographystyle{plainnat}
\bibliographystyle{TRAILbib}


% add here your own bibliography files
\bibliography{./references}

\clearemptydoublepage

%%%%%%%%%%%%%%%%%%%%%%%%%%%%%%%%%%%%%%%%%%%%%%%%%%%%%%%%%%%%%%%%

% here comes the glossary (partially constructed by the commands in
% 0-misc/lists.tex, which was included above)

%\chapter*{Glossary}\markboth{Glossary}{Glossary}
\addcontentsline{toc}{chapter}{Glossary}

%\renewcommand{\imagepath}[1]{\fname{0-misc/images/#1}}
\graphicspath{{./0-misc/images/}}


\section*{Conventions}
The following conventions are used in this thesis for notation
and symbols:
\begin{itemize}
\item A lower case character typeset in boldface, e.g., $\mathbf{x}$,
   represents a column vector.
\item The number of elements in a vector $\mathbf{x}$ is indicated by
   $n_{\mathbf{x}}$.
\item \ldots
\end{itemize}

\section*{List of symbols and notations}
Below follows a list of the most frequently used symbols and notations
in this thesis. Symbols particular to power network applications
are explained only in the relevant chapters.

\vspace*{-0.4cm}\begin{tabbing}
\hspace*{3.5cm}\=\kill\\
$ \mathbf{A}  $ \> system matrices of linear time-invariant models \\
\\[-0.1cm]
$ \mathbf{B}  $ \> input matrices of linear time-invariant models \\
\\[-0.1cm]
$ \mathbf{E}_1 $ \> matrices of mixed-logical dynamic models \\
\end{tabbing}

\vspace*{-0.7cm}
\section*{List of abbreviations}
The following abbreviations are used in this thesis:
\begin{tabbing}
\hspace*{4cm}\=\kill
AVR \> Automatic Voltage Regulator\\
DAE \> Differential-Algebraic Equations\\
FACTS \> Flexible Alternating-Current Transmission System\\
MPC \> Model Predictive Control\\
\ldots \> \ldots
\end{tabbing}
}
\chapter*{Glossary}\markboth{Glossary}{Glossary}
\addcontentsline{toc}{chapter}{Glossary}

%\renewcommand{\imagepath}[1]{\fname{0-misc/images/#1}}
\graphicspath{{./0-misc/images/}}


\section*{Conventions}
The following conventions are used in this thesis for notation
and symbols:
\begin{itemize}
\item A lower case character typeset in boldface, e.g., $\mathbf{x}$,
   represents a column vector.
\item The number of elements in a vector $\mathbf{x}$ is indicated by
   $n_{\mathbf{x}}$.
\item \ldots
\end{itemize}

\section*{List of symbols and notations}
Below follows a list of the most frequently used symbols and notations
in this thesis. Symbols particular to power network applications
are explained only in the relevant chapters.

\vspace*{-0.4cm}\begin{tabbing}
\hspace*{3.5cm}\=\kill\\
$ \mathbf{A}  $ \> system matrices of linear time-invariant models \\
\\[-0.1cm]
$ \mathbf{B}  $ \> input matrices of linear time-invariant models \\
\\[-0.1cm]
$ \mathbf{E}_1 $ \> matrices of mixed-logical dynamic models \\
\end{tabbing}

\vspace*{-0.7cm}
\section*{List of abbreviations}
The following abbreviations are used in this thesis:
\begin{tabbing}
\hspace*{4cm}\=\kill
AVR \> Automatic Voltage Regulator\\
DAE \> Differential-Algebraic Equations\\
FACTS \> Flexible Alternating-Current Transmission System\\
MPC \> Model Predictive Control\\
\ldots \> \ldots
\end{tabbing}

\clearemptydoublepage
%%%%%%%%%%%%%%%%%%%%%%%%%%%%%%%%%%%%%%%%%%%%%%%%%%%%%%%%%%%%%%%%
% here comes the summary in english

%\chapter*{Summary}\markboth{Summary}{Summary}
\addcontentsline{toc}{chapter}{Summary}

Events are getting more popular and more frequent in cities around the world. In the Netherlands in 2017, the number of festivals grew to almost 1000 \footnote{https://www.eventbranche.nl/nieuws/aantal-festivals-groeit-tot-bijna-1000-per-jaar-aantal-bezoeken-daalt-miniem-16483.html}. These events take place in large areas of the city, they have a common topic, they include sub-events (activities), and they have start and end times and lasts from one day to several days. Examples of events are the national holidays, Soul Live Festival and trade exhibitions. City events can easily attract a large number of people. Event stakeholders, such as the event organizers, police, municipalities and other authorities, and crowd managers are concerned with guaranteeing the safety, comfort and general well being of the attendees. To this end, they enforce predefined crowd management measures that are adaptive to the current state of the event environment and of the participating crowd. This state is measured through information about the factors influencing event planning \citep{li2019crowds} and pedestrian behaviour \citep{still2000crowd,tubbs2007egress,abbott2000event, zomer2015managing} for crowd management, such as crowd size, density, mobility, emotion, visitor profile, and location. Conventionally, this information is derived from data provided by stewards (operating on the ground during the event) and sometimes pre-installed sensing infrastructures, such as counting systems, Bluetooth/ Wi-Fi sensors, and video cameras. While effective, these solutions suffer from several issues: they provide little information about sentiments, gender and age distribution, they are expensive, they cannot provide Spatio-temporal information, and they are complex to install and maintain.  

}
\chapter*{Summary}\markboth{Summary}{Summary}
\addcontentsline{toc}{chapter}{Summary}

Events are getting more popular and more frequent in cities around the world. In the Netherlands in 2017, the number of festivals grew to almost 1000 \footnote{https://www.eventbranche.nl/nieuws/aantal-festivals-groeit-tot-bijna-1000-per-jaar-aantal-bezoeken-daalt-miniem-16483.html}. These events take place in large areas of the city, they have a common topic, they include sub-events (activities), and they have start and end times and lasts from one day to several days. Examples of events are the national holidays, Soul Live Festival and trade exhibitions. City events can easily attract a large number of people. Event stakeholders, such as the event organizers, police, municipalities and other authorities, and crowd managers are concerned with guaranteeing the safety, comfort and general well being of the attendees. To this end, they enforce predefined crowd management measures that are adaptive to the current state of the event environment and of the participating crowd. This state is measured through information about the factors influencing event planning \citep{li2019crowds} and pedestrian behaviour \citep{still2000crowd,tubbs2007egress,abbott2000event, zomer2015managing} for crowd management, such as crowd size, density, mobility, emotion, visitor profile, and location. Conventionally, this information is derived from data provided by stewards (operating on the ground during the event) and sometimes pre-installed sensing infrastructures, such as counting systems, Bluetooth/ Wi-Fi sensors, and video cameras. While effective, these solutions suffer from several issues: they provide little information about sentiments, gender and age distribution, they are expensive, they cannot provide Spatio-temporal information, and they are complex to install and maintain.  


\clearemptydoublepage

%%%%%%%%%%%%%%%%%%%%%%%%%%%%%%%%%%%%%%%%%%%%%%%%%%%%%%%%%%%%%%%%
% here comes a summary in dutch

%\chapter*{Samenvatting}\markboth{Samenvatting}{Samenvatting}
\addcontentsline{toc}{chapter}{Samenvatting (Summary in Dutch)}

% Samenvatting should be in dutch, so use dutch hyphenation
\selectlanguage{dutch}

\section*{Social media gebruiken om menigte te karakteriseren in stadsevenementen voor crowd management}

Een samenvatting in het Nederlands zal hier worden gepresenteerd.


\bigskip
\noindent Vincent X.\ Gong

\selectlanguage{american}

%%% Local Variables: 
%%% mode: latex
%%% TeX-master: "../2007-thesis"
%%% End: 
}
\chapter*{Samenvatting}\markboth{Samenvatting}{Samenvatting}
\addcontentsline{toc}{chapter}{Samenvatting (Summary in Dutch)}

% Samenvatting should be in dutch, so use dutch hyphenation
\selectlanguage{dutch}

\section*{Social media gebruiken om menigte te karakteriseren in stadsevenementen voor crowd management}

Een samenvatting in het Nederlands zal hier worden gepresenteerd.


\bigskip
\noindent Vincent X.\ Gong

\selectlanguage{american}

%%% Local Variables: 
%%% mode: latex
%%% TeX-master: "../2007-thesis"
%%% End: 

\clearemptydoublepage
%%%%%%%%%%%%%%%%%%%%%%%%%%%%%%%%%%%%%%%%%%%%%%%%%%%%%%%%%%%%%%%%
% here comes the summary in other languages.
% you can remove it if you do not want it.

%\chapter*{Summary}\markboth{Summary}{Summary}
\addcontentsline{toc}{chapter}{Summary}

Events are getting more popular and more frequent in cities around the world. In the Netherlands in 2017, the number of festivals grew to almost 1000 \footnote{https://www.eventbranche.nl/nieuws/aantal-festivals-groeit-tot-bijna-1000-per-jaar-aantal-bezoeken-daalt-miniem-16483.html}. These events take place in large areas of the city, they have a common topic, they include sub-events (activities), and they have start and end times and lasts from one day to several days. Examples of events are the national holidays, Soul Live Festival and trade exhibitions. City events can easily attract a large number of people. Event stakeholders, such as the event organizers, police, municipalities and other authorities, and crowd managers are concerned with guaranteeing the safety, comfort and general well being of the attendees. To this end, they enforce predefined crowd management measures that are adaptive to the current state of the event environment and of the participating crowd. This state is measured through information about the factors influencing event planning \citep{li2019crowds} and pedestrian behaviour \citep{still2000crowd,tubbs2007egress,abbott2000event, zomer2015managing} for crowd management, such as crowd size, density, mobility, emotion, visitor profile, and location. Conventionally, this information is derived from data provided by stewards (operating on the ground during the event) and sometimes pre-installed sensing infrastructures, such as counting systems, Bluetooth/ Wi-Fi sensors, and video cameras. While effective, these solutions suffer from several issues: they provide little information about sentiments, gender and age distribution, they are expensive, they cannot provide Spatio-temporal information, and they are complex to install and maintain.  

}
% \chapter*{Summary in Chinese}\markboth{Summarycn}{Summarycn}
\addcontentsline{toc}{chapter}{Summary in Chinese}
\begin{CJK*}{UTF8}{gbsn}
A summary in Chinese will be presented here.

摘要。
\end{CJK*}


%\bigskip
%\noindent Vincent X.\ Gong

%%% Local Variables: 
%%% mode: latex
%%% TeX-master: "../2007-thesis"
%%% End: 

% \clearemptydoublepage

%%%%%%%%%%%%%%%%%%%%%%%%%%%%%%%%%%%%%%%%%%%%%%%%%%%%%%%%%%%%%%%%

% and here comes the curriculum vitae of the author

%\chapter*{About the author}\markboth{About the author}{About the author}

\addcontentsline{toc}{chapter}{About the author}

\begin{wrapfigure}{O}{0.34\textwidth}
%\centering
\includegraphics[width=0.34\textwidth, trim=0 0 0 0.8cm]{./0-misc/images/vincent.jpg}
%\caption{\label{fig:frog2}This is a figure caption.}
\end{wrapfigure}

His research interests include social data analysis, crowd behaviour and crowd management in city events.



\vspace*{\fill}
\pagebreak

\vspace*{0.60cm}
\noindent{\Huge \textbf{Publications}}

\section*{Journal papers}

\begin{enumerate}
%  \item \textbf{Gong, Vincent X.}, et al. "Using social media for attendees density estimation in city-scale events." IEEE Access 6 (2018): 36325-36340.
  \item \textbf{Gong, Vincent X.}, Jie Yang, Winnie Daamen, Alessandro Bozzon, Serge P. Hoogendoorn, and Geert-Jan Houben. "Using social media for attendees density estimation in city-scale events." IEEE Access 6 (2018): 36325-36340.
\end{enumerate}

% to get page numbers a multiple of 8 it may be necessary to add some empty pages

%\pagestyle{plain}\fancyhf{}\markboth{}{}
%\newpage{\mbox{}}

%%% Local Variables: 
%%% mode: latex
%%% TeX-master: "../2007-thesis"
%%% End: 
}
\chapter*{About the author}\markboth{About the author}{About the author}

\addcontentsline{toc}{chapter}{About the author}

\begin{wrapfigure}{O}{0.34\textwidth}
%\centering
\includegraphics[width=0.34\textwidth, trim=0 0 0 0.8cm]{./0-misc/images/vincent.jpg}
%\caption{\label{fig:frog2}This is a figure caption.}
\end{wrapfigure}

His research interests include social data analysis, crowd behaviour and crowd management in city events.



\vspace*{\fill}
\pagebreak

\vspace*{0.60cm}
\noindent{\Huge \textbf{Publications}}

\section*{Journal papers}

\begin{enumerate}
%  \item \textbf{Gong, Vincent X.}, et al. "Using social media for attendees density estimation in city-scale events." IEEE Access 6 (2018): 36325-36340.
  \item \textbf{Gong, Vincent X.}, Jie Yang, Winnie Daamen, Alessandro Bozzon, Serge P. Hoogendoorn, and Geert-Jan Houben. "Using social media for attendees density estimation in city-scale events." IEEE Access 6 (2018): 36325-36340.
\end{enumerate}

% to get page numbers a multiple of 8 it may be necessary to add some empty pages

%\pagestyle{plain}\fancyhf{}\markboth{}{}
%\newpage{\mbox{}}

%%% Local Variables: 
%%% mode: latex
%%% TeX-master: "../2007-thesis"
%%% End: 

%\clearemptydoublepage

%%%%%%%%%%%%%%%%%%%%%%%%%%%%%%%%%%%%%%%%%%%%%%%%%%%%%%%%%%%%%%%%
% here comes a list with trail series phd theses. (no need to include
% these if you are not publishing in the trail series.)

% \chapter*{TRAIL Thesis Series}
\addcontentsline{toc}{chapter}{TRAIL Thesis Series publications}
\markboth{TRAIL Thesis Series}{} \markright{TRAIL Thesis Series}

\noindent The following list contains the most recent dissertations in the TRAIL Thesis Series, For a complete overview of more than 250 titles see the TRAIL website: www.rsTRAIL.nl.

\vspace*{0.40cm}{}
\noindent The TRAIL Thesis Series is a series of the Netherlands TRAIL Research School on transport, infrastructure and logistics.

\vspace*{0.40cm}{}
\noindent Nat, C.G.J.M., van der, \textit{A Knowledge-based Concept
Exploration Model for Submarine Design}, T99/1, March 1999, TRAIL
Thesis Series, Delft University Press, The Netherlands

\vspace*{0.25cm}{} 
\noindent Nuttall, A.J.G., \textit{Design Aspects of Multiple Driven Belt Conveyors}, T2007/12, November 2007, TRAIL Thesis Series, The
Netherlands

\vspace*{0.25cm}{} 
\noindent Nederveen, A.A.J., \textit{Ruimtelijke Inpassing van Lijninfrastructuur. Een onderzoek naar de geschiktheid van inspraakreacties voor het ontwerpen van lijninfrastructuur}, T2007/13, December 2007, TRAIL Thesis Series, The
Netherlands

\vspace*{0.25cm}{} 
\noindent Negenborn, R.R., \textit{Multi-Agent Model Predictive Control with
Applications to Power Networks}, T2007/14, December 2007, TRAIL Thesis
Series, The Netherlands


%%% Local Variables: 
%%% mode: latex
%%% TeX-master: "../2007-thesis"
%%% End: 
}
\chapter*{TRAIL Thesis Series}
\addcontentsline{toc}{chapter}{TRAIL Thesis Series publications}
\markboth{TRAIL Thesis Series}{} \markright{TRAIL Thesis Series}

\noindent The following list contains the most recent dissertations in the TRAIL Thesis Series, For a complete overview of more than 250 titles see the TRAIL website: www.rsTRAIL.nl.

\vspace*{0.40cm}{}
\noindent The TRAIL Thesis Series is a series of the Netherlands TRAIL Research School on transport, infrastructure and logistics.

\vspace*{0.40cm}{}
\noindent Nat, C.G.J.M., van der, \textit{A Knowledge-based Concept
Exploration Model for Submarine Design}, T99/1, March 1999, TRAIL
Thesis Series, Delft University Press, The Netherlands

\vspace*{0.25cm}{} 
\noindent Nuttall, A.J.G., \textit{Design Aspects of Multiple Driven Belt Conveyors}, T2007/12, November 2007, TRAIL Thesis Series, The
Netherlands

\vspace*{0.25cm}{} 
\noindent Nederveen, A.A.J., \textit{Ruimtelijke Inpassing van Lijninfrastructuur. Een onderzoek naar de geschiktheid van inspraakreacties voor het ontwerpen van lijninfrastructuur}, T2007/13, December 2007, TRAIL Thesis Series, The
Netherlands

\vspace*{0.25cm}{} 
\noindent Negenborn, R.R., \textit{Multi-Agent Model Predictive Control with
Applications to Power Networks}, T2007/14, December 2007, TRAIL Thesis
Series, The Netherlands


%%% Local Variables: 
%%% mode: latex
%%% TeX-master: "../2007-thesis"
%%% End: 

%\clearemptydoublepage
%%%%%%%%%%%%%%%%%%%%%%%%%%%%%%%%%%%%%%%%%%%%%%%%%%%%%%%%%%%%%%%%

% that's it!
\end{linenumbers}
\end{document}


%%% Local Variables:
%%% mode: latex
%%% TeX-master: t
%%% End:
